\begin{examplebox}
	\textbf{\textcolor{CExample}{例 1（基础）}}

	\textbf{\textcolor{CExample}{题目：}} 在直角三角形中，两直角边长分别为 $3$ 和 $4$，求斜边长。

	\textbf{\textcolor{CExample}{解题步骤：}}
	\begin{description}[style=nextline, leftmargin=2.8em, labelsep=0.5em, itemsep=0.45em]
		\item[\textbf{第一步（代入公式）：}] 由勾股定理得
		      \[
			      c^2 = 3^2 + 4^2 = 9 + 16 = 25.
		      \]
		      \par\textbf{理由：}
		      已知的是两条直角边，所以用 $a^2+b^2=c^2$ 求斜边平方。

		\item[\textbf{第二步（开方求值）：}]
		      \[
			      c = \sqrt{25} = 5.
		      \]
		      \par\textbf{理由：}
		      边长只能取正数。

		\item[\textbf{第三步（作答）：}]
		      \[
			      c=5.
		      \]
		      \par\textbf{理由：}
		      写出最终结果。
	\end{description}

	\textbf{\textcolor{CExample}{关键结论：}} \textbf{斜边长为 $5$。}

	\medskip
	\textbf{\textcolor{CExample}{例 2（稍变形）}}

	\textbf{\textcolor{CExample}{题目：}} 直角三角形中，斜边长 $13$，一条直角边长 $5$，求另一条直角边长。

	\textbf{\textcolor{CExample}{解题步骤：}}
	\begin{description}[style=nextline, leftmargin=2.8em, labelsep=0.5em, itemsep=0.45em]
		\item[\textbf{第一步（变形公式）：}]
		      \[
			      b^2 = c^2 - a^2 = 13^2 - 5^2 = 169 - 25 = 144.
		      \]
		      \par\textbf{理由：}
		      已知斜边和一条直角边，所以用斜边平方减去已知直角边平方，得到另一条直角边的平方。

		\item[\textbf{第二步（开方得边长）：}]
		      \[
			      b = \sqrt{144} = 12.
		      \]
		      \par\textbf{理由：}
		      边长取正值。

		\item[\textbf{第三步（作答）：}]
		      \[
			      b=12.
		      \]
		      \par\textbf{理由：}
		      写出最终结果。
	\end{description}

	\textbf{\textcolor{CExample}{关键结论：}} \textbf{另一条直角边长为 $12$。}

	\medskip
	\textbf{\textcolor{CExample}{例 3（综合应用）}}

	\textbf{\textcolor{CExample}{题目：}} 在直角三角形中，斜边比一条直角边长 $1$，另一条直角边长为 $5$。求这条直角边和斜边的长度。

	\textbf{\textcolor{CExample}{解题步骤：}}
	\begin{description}[style=nextline, leftmargin=2.8em, labelsep=0.5em, itemsep=0.45em]
		\item[\textbf{第一步（设未知数列方程）：}] 设与斜边相差 $1$ 的那条直角边长为 $x$，则斜边长为 $x+1$，另一条直角边长为 $5$。由勾股定理得
		      \[
			      x^2 + 5^2 = (x+1)^2.
		      \]
		      \par\textbf{理由：}
		      $x$ 和 $5$ 是两条直角边，$x+1$ 是斜边，所以直角边平方和等于斜边平方。

		\item[\textbf{第二步（展开并解方程）：}]
		      \[
			      \begin{aligned}
				      x^2 + 25 & = x^2 + 2x + 1 \\
				      25       & = 2x + 1       \\
				      x        & = 12.
			      \end{aligned}
		      \]
		      \par\textbf{理由：}
		      两边同时减去 $x^2$ 并移项。

		\item[\textbf{第三步（求斜边长）：}] 斜边长为
		      \[
			      x+1 = 13.
		      \]
		      \par\textbf{理由：}
		      代入已求得的 $x$ 值。
	\end{description}

	\textbf{\textcolor{CExample}{关键结论：}} \textbf{这条直角边长为 $12$，斜边长为 $13$。}
\end{examplebox}